\selectlanguage{ngerman}%
\frenchspacing%
\begin{ledgroupsized}[r]{120mm}%
\footnotesize%
\pstart%
\noindent%
\textbf{Überlieferung:}%
\pend%
\end{ledgroupsized}%
\begin{ledgroupsized}[r]{114mm}%
\footnotesize%
\pstart%
\parindent -6mm%
\makebox[6mm][l]{\textit{L}}%
Notiz:
LH XXXV 9,\;16 Bl.\;7.
Ein Zettel (6\;x\;4 cm),
Ränder beschnitten.
Eine volle Seite auf Bl.\;\href{\Scan{LH_35_09_16_007v}}{7\,v\textsuperscript{o}}
und drei Zeilen auf Bl.\;\href{\Scan{LH_35_09_16_007r}}{7\,r\textsuperscript{o}},
sonst unbeschrieben;
Textfolge durch einen Kustos (\textit{verte}) am Ende von Bl.\;\href{\Scan{LH_35_09_16_007v}}{7\,v\textsuperscript{o}} festgelegt.
\pend%
\end{ledgroupsized}%
\vspace{5mm}%
\begin{ledgroup}%
\footnotesize%
\pstart%
\noindent%
\textbf{Datierungsgründe:}
\pagelabel{41172_dat}
Die undatierte Notiz N.\;\ref{41172} ist auf einem Träger überliefert, der weder Spuren eines Wasserzeichens noch äußere Merkmale aufweist, die eindeutige Anhaltspunkte für eine Datierung liefern würden.
Das Zettelformat deutet aber auf eine Entstehung der Notiz N.\;\ref{41172} nach Leibnizens Umzug von Paris\indexO{Paris} nach Hannover:\indexO{Hannover}
Kleine und fragmentarische Zettel sind im Nachlass seltener aus der Mainzer\indexO{Mainz} oder Pariser\indexO{Paris} Zeit überliefert.
Demnach ist unwahrscheinlich, dass die Notiz N.\;\ref{41172} vor 1677 verfasst wurde.
\pend%
\pstart%
Im Mittelpunkt des kurzen Textes steht der Gedanke, dass der Widerstand, mit dem sich ein Körper dem Bewegt-Werden widersetzt -- eine Art Trägheit~--, umso mehr wächst, je höher die Geschwindigkeit der Bewegung wird.
Bei dieser Annahme zeigt die Notiz N.\;\ref{41172} enge inhaltliche Verwandschaft mit weiteren Texten der 1670er und frühen 1680er Jahre: zum Beispiel mit der möglicherweise zwischen 1680 und 1684 abgefassten Aufzeichnung \textit{Corpus impingens a penetrabili repercuti potest} im vorliegenden Band (siehe für weitere Einzelheiten die ausführlichen Datierungsgründe zu N.\;\ref{60283}, S.\;\pageref{60283_dat}).
Dieser Befund ist mit der Vermutung vereinbar, dass die Notiz N.\;\ref{41172} erst in der Hannoveraner\indexO{Hannover} Zeit verfasst wurde, d.\,h. nicht vor 1677.
Eine frühere Entstehung ist nicht ausgeschlossen, aber unwahrscheinlich.
Unwahrscheinlich ist ebenfalls, dass die Notiz N.\;\ref{41172} viel später als die besagte Aufzeichnung N.\;\ref{60283} entstand.
Ein eindeutiger Terminus ante quem der Datierung lässt sich jedoch nach heutigem Wissensstand nicht angeben.
\pend%
\end{ledgroup}%
\selectlanguage{latin}%
\frenchspacing%
\vspace{8mm}%
\normalsize%
\pstart%
\noindent%
\lbrack%
\href{\Scan{LH_35_09_16_007v}}%
{7\,v\textsuperscript{o}}%
\rbrack\ %
Examinandum
\edtext{%
an Mensa aliqua%
\indexS{mensa}%
}{%
\lemma{an}%
\Bfootnote{%
\textit{(1)}~corpus %
\textit{(2)}~Mensa aliqua%
~\textit{L}%
}}
vel \textlangle n\textrangle avis%
\indexS{navis} %
tanta celeritate%
\indexS{celeritas motus}
moveri possit,
ut corpus ei incumbens%
\indexS{corpus incumbens altero}
gravitate sua%
\indexS{gravitas corporis}
non sequatur %
sed nonnihil restitet,%
\indexS{corpus restitens}
et in mensa%
\indexS{mensa}
quasi retrorsum cedat,
non ob resistentiam aeris%
\indexS{resistentia aeris}%
\indexS{aer resistens} %
a qua animum abstrahimus,%
\indexS{animus}
\edtext{%
sed ideo
quod eminentiarum%
\indexS{eminentiae elasticae corporum}%
\indexS{elastrum eminentiarum}
(quibus corpus impositum%
\indexS{corpus impositum altero}
teneri videtur)
elastrum cedat%
}{%
\lemma{sed}%
\Bfootnote{%
\textit{(1)}~ob %
\textit{(2)}~ideo quod eminentiarum %
\textit{(a)}~elastrum cedat \textlangle qu\textrangle\ %
\textit{(b)}~(quibus corpus %
\lbrack...\rbrack\ elastrum cedat%
~\textit{L}%
}}
retinentiae corporis%
\indexS{retinentia corporis}
contra magnam celeritatem.%
\indexS{celeritas magna}
Ita puto.
\lbrack%
\href{\Scan{LH_35_09_16_007r}}%
{7\,r\textsuperscript{o}}%
\rbrack\
Ut videmus quaedam corpora
mota
\edtext{%
celeri\lbrack tate\rbrack%
}{%
\lemma{celeri}%
\Bfootnote{%
\textit{L\;ändert Hrsg.}%
}}
frangi,%
\indexS{corpus fractum}%
\indexS{fractura corporis}
dum partes remotiores satis promte sequi non possunt.%
\indexS{pars corporis non sequens motum}
\pend%
\selectlanguage{latin}%
\frenchspacing%
\normalsize%
\newpage%
